\section{Introduction}\label{sec:intro}

\subsection{The question}

Collateralized lending in decentralized finance values collateral at a market price and closes positions by forced sale when that price falls. Forced sales depress prices and can trigger further liquidations \citep{perez2021,qin2021liquidations,lehar2022,chiu2023}. On 12 March 2020, more than \$8 million of ETH collateral in MakerDAO vaults was sold in auctions that cleared at a price of zero \citep{glassnodeblackthursday}.

A different design removes the market price from the loan altogether. Suppose the borrower's collateral is a token that the lender itself issues through a mint-and-burn bonding curve, and that the lender also holds the reserve from which that token is redeemed. The lender then knows a lower bound on the token's redemption value that does not depend on any market: the lowest price on its own curve, the \emph{floor price} $\Pf$. If a loan never exceeds the floor value of the tokens locked behind it, and the floor never falls, the loan never becomes under-collateralized, whatever the market does. No oracle and no liquidation are needed.

The idea is no longer new as a product. Nirvana offered loans at the floor value in 2022, paid out in its own stablecoin \citep{nirvanadocs}. Baseline's current system lends against its floor (``baseline value'') at zero interest, with a one-time origination fee, no expiry and no liquidation threshold \citep{baselinedocs}. Debt-aware bonding curves \vone{} combine a stepped issuance curve with a credit facility of the same kind. What these designs have in common is credit anchored to the issuer's own floor. This paper analyzes the case in which the loan is paid from the reserve, so that the issuer \emph{lends the redemption reserve} to the holders whose tokens that reserve backs. (A loan paid in a token that the issuer creates, as in Nirvana, has a different balance sheet.)

That balance sheet raises three connected questions:
\begin{quote}
\emph{Which credit allocations remain safe under the collateral-withdrawal rights promised to borrowers? How much additional reserve, or restriction on those rights, is needed to lend more? What does this safety imply for borrowers' incentives and willingness to enter?}
\end{quote}
The earlier paper \vone{} gives a coverage argument for stepped curves. We extend that argument to arbitrary non-decreasing schedules and then examine its boundary. A loan book can cover every redemption today and fail after an allowed partial repayment tomorrow. The relevant restriction is therefore a condition on both the amount lent and the collateral that repayments release.

\subsection{Contributions}

\paragraph{Credit capacity determined by withdrawal rights.}
\Cref{sec:credit} fixes the curve and characterizes coverage after every permitted repayment and collateral release (\cref{cr:kernel}). The current-state condition $D\le R(L)$ without surplus is insufficient for that stronger requirement. With unrestricted proportional release and no committed reserve beyond curve backing, every individual loan must satisfy the floor bound, even when loan limits are heterogeneous or depend on the state (\cref{cr:buffer}). With a buffer $W$, a sorted debt-per-collateral envelope $H$ gives the exact requirement
\[
W\ge\max_{0\le q\le L}\{H(q)-R(q)\}.
\]
Equal debt ratios weakly minimize that requirement and attain the greatest aggregate credit for a fixed buffer (\cref{cr:capacity}). The envelope is classical fractional knapsack, equivalently an upper integrated-quantile curve; the contribution is its connection to the states made attainable by the repayment contract. The requirement has a second meaning, which is the result we consider most consequential: under stated zero-fee coalition assumptions it equals the most that borrowers can gain by repaying in part, redeeming the released collateral and abandoning the rest, and at the maximum admitted credit that strategy can exhaust a reserve funded by others (\cref{cr:extraction}). Payment safety and the interest of whoever funds the safety are different things. The appendix gives an exact finite algorithm and shows why checking indivisible exit rights instead contains a subset-selection problem (\cref{cr:complexity}).

\paragraph{A different contract, with different costs.}
A loan-local cap based on $R(\ell)$, the reserve value of the entire lock, can permit more credit than the floor value if partial releases preserve the same nonlinear bound (\cref{cr:nonlinearsafety}). It requires a stronger reshape guard and gives less collateral back for a repayment at the loan limit. Within this zero-buffer alternative, pooling can increase capacity by keeping collateral subject to a joint release condition. Arbitrary proportional splitting into independently repayable loans recovers the floor bound (\cref{cr:splitting}); small locks within the flat part of a stepped curve obtain no extra capacity. These are analytically specified alternatives. The supplied Solidity implementation remains the floor-cap contract.

\paragraph{The foundation: a complete floor-cap safety theorem.}
\Cref{sec:model} states three invariant conditions: locked tokens are in custody and cannot be redeemed (J1); each loan is bounded by a fixed fraction of its collateral's floor value (J2); and cash plus debt covers the curve reserve of the whole supply (J3). \Cref{thm:safety} proves their preservation over the full operation set and derives payability of every free-token redemption and admissible loan, without assuming any repayment. Curve changes need a non-decreasing floor and adequate funding. The argument is elementary, independent of discreteness, and formalized in Lean~4. Its contribution is a complete specification whose operations and guards can be compared with an implementation.

\paragraph{Incentives and the participation constraint.}
For that floor-cap contract, \cref{sec:incentives} analyzes repayment as the exercise of a perpetual option with no carry, following the stock-loan interpretation of \citet{xiazhou2007}. Deferring repayment until the tokens are used is weakly optimal, and borrowing benefits a holder with a sufficiently long horizon. A cap below the net-floor ratio guarantees a collateral-funded close and removes the cash advantage of borrowing and abandoning over redeeming. At the floor, with arrivals given, the optimal plans aggregate into a stationary loan book and a floor-growth rate (\cref{thm:eq}). But when growth is below the discount rate and the cap is below the net floor, the mean entry value is negative (\cref{thm:entry}). The mechanism's cash flows then cannot explain positive voluntary entry. In a numerical example the missing benefit is $1.16\%$ of the floor price per token and year. A refundable-capital benchmark shows what a service benefit would have to compensate and why a deposit offering the same service is an essential comparison (\cref{prop:deposit-benchmark}). It is not a model of endogenous adoption.

\paragraph{Redistribution, extraction and executable evidence.}
Coverage does not prevent a reconfiguration from moving exit value between holders (\cref{sec:reconf}). \Cref{sec:extraction} derives exact sandwich and repricing profits and states the limits of the corresponding guards. An integer reference model, a floor-cap Solidity implementation and differential traces support the original mechanism. Separate exact-rational checks support the new credit-capacity calculations, including adverse repayment paths and changes that deliberately omit a required guard (\cref{sec:evidence}). The formal proof, testing and analytic results have distinct scopes.

\paragraph{Evidence from a deployed design.}
\Cref{sec:empirics} examines \ELogs{} Baseline event logs on Ethereum over \EDays{} days. In the pool with the most accounts, \EBShare\% of circulating tokens are locked; plain repayments are followed within about a week by a sale by the same account in \ERepaySoldPct\% of cases, or \ERestPct\% after excluding the two most active accounts. By amount, \ELeverAmtPct\% of new gross debt finances token purchases. A new reconciliation of fee-distribution events separates allocations to protocol, creators and staking from separately recorded claims, and matches origination and swap fees conservatively. This documents a reward channel absent from the no-payout benchmark. It does not identify borrowers' earnings, beliefs or returns. All empirical findings are descriptive.

\subsection{Relation to the earlier paper}

This paper follows \vone, which introduced debt-aware discrete bonding curves and of which the present author is a co-author. It is written to stand alone. It generalizes the core result of that paper and adds the analysis summarized above. \Cref{app:errata} lists where its statements differ from the earlier ones. The difference that matters most for anyone building on the earlier text is the reshape guard: a reconfiguration must also keep the floor from falling, (G1), because solvency depends on it (\cref{ex:reconf}).

\subsection{Scope}

The floor-cap safety theorem and the credit-capacity extensions concern \emph{issuance-native} collateral only: tokens minted by the curve and redeemable from its reserve. They are stated in units of the reserve asset. A floor that never falls in units of a reserve token says nothing about the value of that token. The floor is a primary-market redemption price before fees; a secondary market can trade below it (\cref{sec:secondary}). We do not model the demand for the token and make no claim about whether such a token is a good investment. The credit-capacity characterization fixes the curve and specifies the allowed repayment rights; its dynamic extensions require their own guards. The incentive results return to the floor-cap contract and compare the plans of a single price-taking agent under stated assumptions; the aggregate result covers a market at the floor with given arrivals and not the interdependent exit decisions above it. The empirical section is descriptive and concerns one protocol. We do not claim that any deployed protocol satisfies all of the model's assumptions; \cref{sec:conformance} states what that would require.
